\chapter{Testing for Goodness of Fit}
\label{ch:gof}

We close with a question that has run quietly beneath the whole second half of
the book: does a proposed probability model actually \emph{fit} the data? When we
modeled bomb hits as Poisson or pea types by Mendel's ratios, we compared
observed and predicted frequencies informally. This chapter makes the comparison
rigorous through the likelihood ratio test and Pearson's chi-square test.

\section{Is There a ``Best'' Test?}

First, some vocabulary. A hypothesis fixing a parameter exactly, like
$\theta=\theta_0$, is \emph{simple}; one allowing a range, like $\theta>\theta_0$,
is \emph{composite}. A test of $H_0:\theta=\theta_0$ against $H_1:\theta\neq\theta_0$
is \emph{two-sided}; a test against $H_1:\theta>\theta_0$ is \emph{one-sided}.
Recall the two error probabilities, $\alpha$ (type I) and $\beta$ (type II); the
\emph{power} of a test is $1-\beta$, its chance of correctly rejecting a false
$H_0$. The usual strategy fixes $\alpha$ and seeks the largest possible power. A
test achieving the most power among all level-$\alpha$ tests is \emph{most
powerful}.

Most powerful tests are hard to find and often do not exist. But in the simplest
case they do, and they take a memorable form.

\begin{theorem}[Neyman--Pearson lemma]
To test the simple $H_0:\theta=\theta_0$ against the simple $H_1:\theta=\theta_1$,
form the likelihood ratio
\[
  T=\frac{L(\theta_1)}{L(\theta_0)}=\frac{\prod_i f(x_i;\theta_1)}{\prod_i f(x_i;\theta_0)},
\]
and reject $H_0$ when $T>k$. Choosing $k$ so that $\Prob_{\theta_0}(T>k)=\alpha$
yields the most powerful level-$\alpha$ test.
\end{theorem}

In words: when both hypotheses are simple and there are no nuisance parameters,
the optimal test compares how well the two parameter values explain the data,
through a ratio of likelihoods.

\section{The Likelihood Ratio Test}

Real problems usually involve composite hypotheses or nuisance parameters, where
the Neyman--Pearson lemma does not directly apply. The \emph{likelihood ratio
test} generalizes its idea. To test $H_0:\theta\in\Theta_0$ against
$H_1:\theta\notin\Theta_0$, the test statistic is
\[
  \lambda=2\ln\!\left(\frac{\max_{\theta\in\Theta}L(\theta)}{\max_{\theta\in\Theta_0}L(\theta)}\right)
   =2\ln\!\left(\frac{L(\hat\theta)}{L(\hat\theta_0)}\right),
\]
the (twice log) ratio of the best fit overall to the best fit allowed under
$H_0$. Large $\lambda$ means $H_0$ explains the data far worse than the
unrestricted model, and counts against $H_0$. Its calibration is a theorem.

\begin{theorem}
Under $H_0$, as $n\to\infty$ the statistic $\lambda$ converges to a chi-square
distribution whose degrees of freedom equal the difference in dimension between
$\Theta$ and $\Theta_0$. The $p$-value is the upper-tail probability
$\Prob(\chi^2_{r-q}>\lambda)$.
\end{theorem}

The chi-square family (Figure~\ref{fig:chisq}) thus governs goodness-of-fit
testing, just as the normal and $t$ families governed estimation.

\begin{figure}[ht]
  \centering
  \includegraphics[width=0.55\linewidth]{chisq}
  \caption{Chi-square densities. The degrees of freedom—the difference in
  dimension between the full and restricted parameter spaces—determine which
  curve calibrates the test.}
  \label{fig:chisq}
\end{figure}

\section{The Multinomial Distribution}

Goodness-of-fit data are counts in categories, which follow a \emph{multinomial}
distribution—the multivariate cousin of the binomial. Drawing $n$ times
independently from $k$ categories with probabilities $p=(p_1,\dots,p_k)$, and
letting $X_j$ count category $j$, gives $X\sim\Mult(n,p)$ with
\[
  \Prob(X=x)=\binom{n}{x_1\cdots x_k}p_1^{x_1}\cdots p_k^{x_k}.
\]
Each $X_j$ is marginally $\Bin(n,p_j)$, so $\E[X_j]=np_j$.

\begin{example}[The multinomial MLE]
Show that the MLE of $p$ for multinomial data is $\hat p_j=X_j/n$.

\solution Ignoring a constant, the loglikelihood is
$\ell(p)=\sum_j X_j\ln p_j$, to be maximized subject to $\sum_j p_j=1$. Using a
Lagrange multiplier $\lambda$, maximize
$\sum_j X_j\ln p_j+\lambda(\sum_j p_j-1)$. Setting
$\partial/\partial p_j=\tfrac{X_j}{p_j}+\lambda=0$ gives $\hat p_j=-X_j/\lambda$,
and the constraint $\sum_j\hat p_j=1$ forces $\lambda=-n$, so
$\hat p_j=X_j/n$—each category's observed proportion.
\end{example}

\section{Pearson's Chi-Square Test}

To test a fully specified $H_0:p=p_0$ against $H_1:p\neq p_0$, Pearson's
chi-square statistic compares observed counts $X_j$ to the counts $E_j=np_j$
expected under $H_0$:
\[
  T=\sum_{j=1}^k\frac{(X_j-E_j)^2}{E_j}.
\]
Under $H_0$, $T$ converges to $\chi^2_{k-1}$ (one degree of freedom is lost to the
constraint $\sum X_j=n$), and the $p$-value is $\Prob(\chi^2_{k-1}>t)$. Pearson's
statistic is in fact a second-order Taylor approximation to the likelihood ratio
statistic $\lambda$—a computational shortcut from $1900$, when logarithms were
laborious—so the two give nearly identical answers.

\begin{example}[Mendel's peas]
Mendel's theory predicts that round-yellow, wrinkled-yellow, round-green, and
wrinkled-green peas occur in proportions
$p_0=(\tfrac{9}{16},\tfrac{3}{16},\tfrac{3}{16},\tfrac{1}{16})$. In $n=556$ plants
he observed $X=(315,101,108,32)$. Test $H_0:p=p_0$.

\solution The expected counts are $np_0=(312.75,104.25,104.25,34.75)$
(Figure~\ref{fig:mendel}). Pearson's statistic is
\[
  t=\frac{(315-312.75)^2}{312.75}+\frac{(101-104.25)^2}{104.25}+\frac{(108-104.25)^2}{104.25}+\frac{(32-34.75)^2}{34.75}=0.47.
\]
With $k=4$ categories there are $k-1=3$ degrees of freedom, so the $p$-value is
$\Prob(\chi^2_3>0.47)=0.93$—no evidence against the theory. The likelihood ratio
statistic gives almost the same value, $\lambda=0.48$, with $p$-value $0.92$. By
either test, Mendel's data are strikingly consistent with his $9:3:3:1$ law.
\end{example}

\begin{figure}[ht]
  \centering
  \includegraphics[width=0.55\linewidth]{mendel}
  \caption{Mendel's observed pea counts against the $9:3:3:1$ prediction. The
  near-perfect agreement is reflected in a tiny chi-square statistic.}
  \label{fig:mendel}
\end{figure}

\section{Goodness of Fit with Estimated Parameters}

Often the model under test is not fully specified—we hypothesize a Poisson
distribution but must estimate its mean from the data. Each parameter estimated
costs one further degree of freedom.

\begin{example}[Are accidents Poisson?]
The number of accidents per week at an intersection over $n=50$ weeks was:
\[
\begin{array}{c|cccc}
x & 0 & 1 & 2 & \ge 3\\\hline
\text{frequency} & 32 & 12 & 6 & 0
\end{array}
\]
Test whether the weekly count is Poisson.

\solution Under $H_0$ the count is $\Pois(\mu)$ with $\mu$ unknown. Its MLE is the
total accidents over total weeks, $\hat\mu=24/50=0.48$. Pooling into the three
categories $X=0$, $X=1$, $X\ge 2$, the estimated probabilities are
\[
  \hat p_1=e^{-0.48}=0.619,\quad \hat p_2=0.48e^{-0.48}=0.297,\quad \hat p_3=1-\hat p_1-\hat p_2=0.084,
\]
giving expected counts $E_1=30.95$, $E_2=14.85$, $E_3=4.20$. Then
\[
  t=\frac{(32-30.95)^2}{30.95}+\frac{(12-14.85)^2}{14.85}+\frac{(6-4.20)^2}{4.20}=1.354.
\]
With $k=3$ categories we lose one degree of freedom to the constraint
$\sum X_j=n$ and \emph{another} to estimating $\mu$, leaving $k-2=1$. The
$p$-value is $\Prob(\chi^2_1>1.354)=0.24$, so we do not reject $H_0$: the data
are consistent with a Poisson model for the accident counts.
\end{example}

\bigskip
This brings our tour full circle. We began by describing data, built the
probability theory to model it, and developed the inferential tools—estimation,
regression, and testing—to draw conclusions from it. Goodness-of-fit testing
closes the loop, letting us ask of any model the most basic question of all:
does it fit?
